
% 附录起始位置
\appendix

\chapter{九阴真经原本}

\section{总纲（核心心法）}

\subsection{梵文总纲（后由郭靖、黄蓉译解）：}

原版《九阴真经》的总纲以梵文写成，蕴含武学至高哲理，强调“阴阳互济、刚柔并重”，是化解经中武功戾气的关键。

关键理念：“天之道，损有余而补不足”（出自《道德经》），主张内力修炼需顺应自然，调和阴阳。


\chapter{黯然销魂掌秘籍}

\section{黯然销魂掌的来历}

\subsection{创招背景：}

杨过与小龙女分离十六年，饱受相思之苦，内心极度悲怆。

在南海之滨练剑时，结合毕生所学，创出这套以情驭劲的掌法。

\subsection{武学根基：}

融合《九阴真经》的刚柔并济、古墓派的轻灵迅捷、欧阳锋的蛤蟆功爆发力、黄药师的弹指神通巧劲，以及洪七公的打狗棒法变化。

\section{黯然销魂掌的十七式}

\begin{table}[!h]
    \caption{黯然销魂掌招式解析}
    \begin{tabular}{l l}
        \toprule
        \textbf{招式} & \textbf{要领} \\
        \midrule
        拖泥带水 & 缠绵悱恻，掌力如浪潮般连绵不绝 \\
        孤形只影 & 身形飘忽，掌劲忽左忽右，难以捉摸 \\
        心惊肉跳 & 以极快掌法扰乱对手心神，使其胆寒 \\
        徘徊空谷 & 掌力回荡，如幽谷回声，劲力反复叠加 \\
        力不从心 & 看似无力，实则暗藏后劲，后发制人 \\
        行尸走肉 & 身法诡异，如僵尸般飘忽不定 \\
        庸人自扰 & 掌势混乱无序，却暗合武学至理 \\
        饮恨吞声 & 掌力内敛，蓄势待发，一击必杀 \\
        六神不安 & 掌风笼罩对手全身，使其难以招架 \\
        穷途末路 & 绝境反击，掌力爆发至极限 \\
        面无人色 & 掌风阴寒，令对手如坠冰窟 \\
        想入非非 & 虚招惑敌，实则暗藏杀机 \\
        呆若木鸡 & 静立不动，却蕴含极强防御反击之势 \\
        神不守舍 & 掌法飘忽，如鬼魅般难以预测 \\
        魂不附体 & 掌劲透体，直击对手经脉 \\
        倒行逆施 & 逆反常理，出招角度刁钻 \\
        黯然销魂 & 终极杀招，需极度悲痛心境才能施展 \\
        \bottomrule
    \end{tabular}
\end{table}

\section{黯然销魂掌的特点}